
\part{Organisation des meta-données}
\label{part:meta}
   Le terme de \textbf{meta-données} désigne les 
<<données sur les données>>, c'est à dire toutes
les informations sur les conditions de la mesure
et du traitement de celle-ci.

   Dans le cadre d'AIDA, on considère (principe d'auto-documentation)
que chaque vecteur contient l'information sur les 
conditions de sa propre production (attributs \textit{name},
\textit{unit}, \textit{origine},
\textit{fct\_origine}, \ldots). 

   Il reste donc à décrire : 
\begin{itemize}
  \item les informations propres à la configuration d'AIDA sur une machine
      particulière (liste et caractéristiques des plates-formes connues,
      liste des séries de mesures disponibles)
  \item des informations propres à l'architecture d'AIDA (par exemple
   la façon de regrouper les variables en familles).
  \item les informations relatives à une série de mesure (par exemple, 
      la description d'un vol, son découpage en paliers, \ldots) 
  \item les informations relatives aux campagnes de mesures
\end{itemize}

Il est important que toutes ces informations, même
celles qui évoluent peu, soient contenues dans des fichiers de
type texte, avec un format simple, car cette méthode est la
seule qui garantisse que les différents programmes, écrits dans
différents langages (C, Python, ...) partagent la même
information, et travaillent de façon cohérente.

Ces fichiers sont regroupés dans les répertoires : \\
\-\hspace{1cm} \path{${AIDA_META}} et \path{${AIDA_HOME}/config}. 

A l'exception notable du fichier donnant les détails
de chaque "acquisition" (ou période de mesure), tous
ces fichiers sont organisés de la même façon,
dans le format INI. Ces fichiers sont divisés en paragraphes débutant par un
nom écrit entre
crochets (par ex. \texttt{[merlin]}), suivi par des
lignes du type \texttt{attribut~=~valeur}. Les 
déclarations multi-lignes sont interdites, et les 
commentaires sont introduits par le caractère "\#"(dièse).

Ces fichiers n'ayant pas tous le même rôle, ils n'évolueront pas
de la même manière. Par exemple, ceux qui correspondent 
à l'organisation d'AIDA seront identiques sur
différents ordinateurs, et ne changeront que très
rarement. Ceux qui décrivent la liste des séries de
mesures disponibles évolueront rapidement, et pourront être
différents d'un site à l'autre.

\section{Description des plates-formes}
\label{sec:platef}
   La liste des plates-formes de mesure connues d'AIDA se trouve dans
le fichier \texttt{plates-formes}, dans le répertoire
\path{${AIDA_HOME}/profils/${AIDA_PROFIL}/}. Ce
fichier sert principalement à établir la relation entre le nom standard d'une
plate forme et l'abréviation qui permet de construire les
identificateurs uniques de chaque série de mesure 
(cf section \ref{sec:AcqId}), mais il contient aussi une
description un peu plus détaillée, qui est utilisée pour
renseigner les fichiers de documentation produits par
certains programmes de diffusion de données.

\subsection{Format}
  C'est un fichier au format "INI", divisé en paragraphes. Chaque
  paragraphe correspond à une plate-forme, et se présente comme suit :
\begin{description}
  \item[[Nom de la plate-forme]] tel qu'il est utilisé partout dans AIDA.
  \item[nom\_long] description en une ligne de texte, en français.
  \item[abrev] chaîne de deux caractères.
  \item[aliases] liste d'autres noms qui désignent parfois la plate-forme,
    séparés par des virgules.
  \item[nom\_serie\_mesure] nom d'une période de mesure dans la liste :
    Journee, Vol, Lacher, Lacher\_en\_vol, Drone  (la périodicité "Acquisition"
    n'est plus supportée et est équivalente à Journee)
  \item[duree\_vecteur] indique quelle est la durée d'un vecteur dans AIDA pour
    cette plate-forme. Les valeurs utilisées sont : "journee", "acquisition".
    Cette information est utilisée par les outils d'importation de données
    dans AIDA vers le format vecteur.
  \item[position] pour les plate-formes fixes, indique la position selon le
    format "latitude/longitude/altitude". Pour indiquer explicitement qu'une
    plate-forme n'est pas fixe, indiquer "mobile".\\
    Les coordonnées latitude et longitude sont écrites en degrés décimaux,
    positifs respectivement vers le Nord et l'Est.
\end{description}

  De manière optionnelle, et selon le type de plate-forme, des indicatifs
et/ou immatriculations peuvent être indiqués. Les syntaxes de ces champs sont :
\begin{description}
  \item[ICAO\_indicator] pour les plate-formes situées sur un aéroport,
    indicatif OACI de celui-ci.
  \item[WMO\_station\_index\_number] pour les stations synoptiques,
    indicatif OMM.
  \item[call-sign] pour les avions, immatriculation.
\end{description}

  Enfin, d'autres informations optionnelles peuvent être mentionnées :
\begin{description}
  \item[responsable] nom du reponsable de la plate-forme
\end{description}

\section{Description des acquisitions}
Pour chaque acquisition, il existe un fichier contenant toutes
ses méta-données.

Il est utilisable par les logiciels de visualisation et traitement,
ou simplement lisible par l'utilisateur. Il facilite le travail
de dépouillement des données.

Le format de ce fichier est décrit dans ce chapitre.

\subsection{Nom du fichier}

Pour une acquisition donnée le fichier ``inf'' s'appelle :
\texttt{\$\{Acq\_Id\}.inf}

 Acq\_Id est l'identificateur d'acquisition AIDA. Par exemple,
\path{fa110015} pour la quinzième acquisition de l'année 2011
sur le site du Fauga.
    Il n'y a aucune correspondance entre le numéro et la date de mesure.


\subsection{Stockage}

Tous les fichiers sont rangés dans le répertoire : 

\texttt{\$AIDA\_META\//INF}


\subsection{Format}

C'est un fichier de type texte, constitué d'une série de paragraphes,
eux-mêmes constitués d'une série d'égalités "champ = valeur". 

Le début d'un paragraphe est défini par la présence de son
nom en majuscules et entre crochets suivi d'un retour à la ligne.

A l'intérieur d'un paragraphe, la fin de ligne est marquée
par un point virgule:~\texttt{;}.

A l'intérieur de la valeur les informations sont séparées par un caractère:
\texttt{/}. Ces informations ne doivent pas contenir ce caractère \texttt{/}.

Les lignes vides, les espaces, et les tabulations sont permis.

Les lignes de commentaires doivent commencer par un caractère
dièse: \texttt{\#};

Les champs ne doivent pas contenir d'accents.

Les valeurs peuvent contenir des accents, à condition d'utiliser
le codage iso8859-1 (ou latin1).

Le fichier doit se terminer par \texttt{{[}FIN{]}}.

La liste des paragraphes ainsi que leur caractère "obligatoire"
ou "optionnel" est présentée dans le tableau ci-dessous.

\begin{center}

\begin{tabular}{|l|l|}
\hline 
Nom du Paragraphe& Caractère\\
\hline 
\hline 
\textbf{FORMAT}& \textbf{obligatoire}\\
\hline 
\textbf{IDENTIFICATION}& \textbf{obligatoire}\\
\hline 
\textbf{PARAMETRES\_ACQUIS}&\textbf{obligatoire}\\
\hline 
\textbf{ACQUISITION}& \textbf{obligatoire}\\
\hline 
HORAIRES\_UTC & optionnel\\
\hline 
INS& optionnel\\
\hline 
SEQUENCES& optionnel\\
\hline 
PERSONNEL& optionnel\\
\hline 
POINTS\_CARACTERISTIQUES& optionnel\\
\hline 
METEO\_DEBUT& optionnel\\
\hline  
METEO\_FIN& optionnel\\
\hline 
ETAT\_CAPTEURS& optionnel\\
\hline 
\end{tabular}
\end{center}


\subsubsection{{[}FORMAT{]}}

\label{parag-FORMAT}
Ce paragraphe décrit les formats utilisés dans les valeurs.

\ttfamily
\saute
\noindent\begin{tab2col}
\hline 
\textnormal{Champ}& \textnormal{Valeur}\\
\hline 
\hline 
format\_coord &
C (D)DDdMM.X(X) (C=lettre (N, E, S ou W), (D)DD=degrés 2 ou 3 chiffres, MM=minutes,
X(X)=fraction décimale de minutes;\\
\hline 
format\_heure&
HHMMSS (HH=heures, MM=minutes, SS=secondes);\\
\hline 
format\_date&
JJ/MM/AAAA (JJ=jour, MM=mois, AAAA=année, le tout en chiffres);\\
\hline 
format\_duree&
HH.D (HH=heure, D=dixièmes d'heure);\\
\hline 
\end{tab2col}
\saute

\normalfont

\subsubsection{{[}IDENTIFICATION{]}}
\label{parag-IDENTIFICATION}


\saute
\noindent\begin{tab4col}
\hline 
Champ& Caractère& Format valeur& Description \\
\hline 
\hline 
\texttt{plate-forme}& \textbf{obligatoire} & \textit{texte;}&
nom de la plate-forme en toutes lettres \\
\hline 
\texttt{acq\_id}& \textbf{obligatoire} & \textit{texte;} & identificateur d'acquisition AIDA\\
\hline 
\texttt{date}& \textbf{obligatoire} & \textit{JJ/MM/AAAA;} &
JJ= jour, MM= mois, AAAA= année (format décrit dans le paragraphe FORMAT)\\
\hline 
\texttt{campagne}& \textbf{obligatoire} & \textit{texte;}& nom de la campagne\\
\hline 
\texttt{theme}& optionnel& \textit{texte;}& thème scientifique de l'acquisition\\
\hline 
\texttt{lieu\_depart}& optionnel& \textit{texte;}&
lieu de départ de la plate-forme de mesure\\
\hline 
\texttt{lieu\_arrivee}& optionnel& \textit{texte;}&
lieu d'arrivée de la plate-forme de mesure\\
\hline 
\end{tab4col}
\saute

 Le nom de la plate-forme doit être connu d'AIDA par le fichier
\texttt{plate-forme}

Le champ \texttt{campagne} doit contenir une seule valeur.
Les acquisitions faites
hors de la période opérationnelle de la campagne ont comme nom de
campagne ESSAI, CALIB, ou TRANSIT.

Le champ \texttt{theme} doit expliciter plus librement le but de 
l'acquisition (ex : essais réfractomètre, vol type P1-flux, 
transit Brétigny-Dole, ....).

\subsubsection{{[}HORAIRES\_UTC{]}}

Ce paragraphe est utilisé dans le cas où la plate-forme de mesure est "mobile"
(avion, ballon, drone à ailes...)

\newcommand{\formathoraire}{\textit{hhmmss} ou\newline \textit{AAAA}-\textit{MM}-\textit{JJ}T\textit{hh}:\textit{mm}:\textit{ss}}
\saute
\noindent\begin{tab3col}
\hline 
Champ& Format valeur& Description\\
\hline 
\hline 
\texttt{depart\_parking}& \formathoraire & heure UTC de départ du parking\\
\hline 
\texttt{decollage}& \formathoraire & heure UTC de décollage\\
\hline 
\texttt{atterrissage}& \formathoraire & heure UTC d'atterrissage\\
\hline 
\texttt{retour\_parking}& \formathoraire & heure UTC de retour au parking\\
\hline 
\texttt{duree\_reglementaire}& \textit{hh.d}& durée réglementaire du vol\\
\hline 
\end{tab3col}
\saute


L'existence de deux formats impose aux programmes qui utilisent
ces données de savoir décoder les deux. La distinction se fait sur la
longueur du champ. La forme courte n'indique qu'une heure, et est donc ambigüe
dans les cas où l'acquisition s'étend sur plusieurs jours. Les différents cas
possibles sont :
\begin{itemize}
 \item utilisation d'un champ \textit{hh} >=24. Dans ce cas, l
 \item utilisation d'un champ \textit{hh} <24, mais avec une fin 
antérieure au début (par exemple \texttt{decollage = 223344} et 
\texttt{atterrissage = 001122})
\end{itemize}


Pour l'écriture de 
La forme longue du champ est recommandée, car elle n'est pas ambigüe,
et suit un format reconnu (ISO8601) par de nombreux langages et 
bibliothèques informatiques.

On reconnait que, lorsque les données sont saisies manuellement, 
la forme courte est plus concise et rapide. Il est cependant recommandé 
d'appliquer rapidement une normalisation après


\subsubsection{{[}ACQUISITION{]}}

\saute
\noindent\begin{tab4col}
\hline 
Champ& Caractère & Format valeur& Description \\
\hline 
\hline 
\texttt{systeme\_acquisition}& \textbf{obligatoire} & \textit{texte;}& nom du système d'acquisition\\
\hline 
\texttt{debut\_acquisition}& \textbf{obligatoire} & \textit{JJ/MM/AAAA/hhmmss;}& 
date et heure de début d'acquisition \\
\hline 
\texttt{fin\_acquisition}& \textbf{obligatoire} & \textit{JJ/MM/AAAA/hhmmss;}&
date et heure de fin d'acquisition \\
\hline 
\texttt{media}& \textbf{obligatoire} & \textit{texte;}& 
support de stockage du système d'acquisition\\
\hline 
\texttt{autre\_acquisition}& optionnel & \textit{texte/texte/...;}&
nom d'autres acquisitions sur la m\^eme plate-forme\\
\hline 
\texttt{commentaire}& optionnel & \textit{texte;}& commentaire\\
\hline 
\end{tab4col}
\saute

\noindent\texttt{media} peut prendre les valeurs floppy, dat, exabyte, hd,
don

\texttt{systeme\_acquisition} peut prendre les valeurs :
\begin{itemize}
\item RESUM, CL, PP, PPI pour l'ARAT
\item MC, MC95  pour le Merlin IV
\item EDISOL, ... pour les stations sols 
\item 4M-BATEAU, .... pour les acquisitions 4M  sur bateau
\item STAR, ... pour les radio-sondages
\end{itemize}

\subsubsection{{[}INS{]}}

Ce paragraphe donne les principales informations relatives à
la centrale à inertie de la plate-forme.

\saute
\noindent\begin{tab3col}
\hline 
Champ& Format valeur& Description\\
\hline 
\hline 
\texttt{ramp\_depart}& \textit{Nxxdxx.x/Wxxxdxx.x/xxxxft;}&
coordonnées du point de départ\\
\hline 
\texttt{type\_aln}&
\textit{xxx}~\texttt{\textit{s}}\textit{;} ou \textit{xxx}~\texttt{\textit{mn}}\textit{;}&
durée de l'alignement en secondes ou minutes\\
\hline 
\texttt{ramp\_arrivee}&
\textit{Nxxdxx.x/Wxxxdxx.x/xxxxft;}& coordonnées du point d'arrivée\\
\hline 
\texttt{pp\_arrivee}& \textit{Nxxdxx.x/Wxxxdxx.x/xxxxft;}&
coordonnées réelles à l'arrivée\\
\hline 
\end{tab3col}
\saute

\subsubsection{{[}POINTS\_CARACTERISTIQUES{]}}

Ce paragraphe décrit les différents points 
géographiques intéressant l'acquisition.

\saute
\noindent\begin{tab3col}
\hline 
Champ& Format valeur& Description\\
\hline 
\texttt{wpt}& \textit{nom/Nxxdxx.x/Wxxxdxx.x;}&
abréviation et coordonnées du point \\
\hline 
\texttt{radar}& \textit{nom/Nxxdxx.x/Wxxxdxx.x;}&
abréviation et coordonnées du point \\
\hline 
\texttt{balise}& \textit{nom/Nxxdxx.x/Wxxxdxx.x;}&
abréviation et coordonnées du point \\
\hline 
\texttt{bouee}& \textit{nom/Nxxdxx.x/Wxxxdxx.x;}&
abréviation et coordonnées du point \\
\hline 
\texttt{station}& \textit{nom/Nxxdxx.x/Wxxxdxx.x;}&
abréviation et coordonnées du point \\
\hline 
\texttt{ville}& \textit{nom/Nxxdxx.x/Wxxxdxx.x;}&
abréviation et coordonnées du point \\
\hline 
\end{tab3col}
\saute

\subsubsection{{[}PERSONNEL{]}}

\saute
\noindent\begin{tab3col}
\hline 
Champ&
Format valeur&
Description\\
\hline 
\hline 
\texttt{pilote}& \textit{Initiale du prénom. Nom;}&
Nom du pilote de la plate-forme (avion)\\
\hline 
\texttt{copilote}& idem&
Nom du copilote de la plate-forme\\
\hline 
\texttt{operateur}& idem&
Nom de l'opérateur du système d'acquisition\\
\hline 
\texttt{scientifique}& idem&
Nom du scientifique de la mission\\
\hline 
\texttt{directeur\_de\_mission}& idem&
Nom du directeur de la mission\\
\hline 
\texttt{invite}& idem&
Nom de l'invité \\
\hline 
\texttt{commandant\_de\_bord}& idem&
Nom du commandant de bord (bateau)\\
\hline 
\end{tab3col}
\saute

Les champs peuvent être répétés si plusieurs personnes
correspondent à la même désignation (ex: plusieurs 
scientifiques).


\subsubsection{{[}METEO\_DEBUT{]}}

Ce paragraphe décrit les conditions météorologiques 
en début d'acquisition.

\saute
\noindent\begin{tab3col}
\hline 
Champ& Format valeur& Description\\
\hline 
\hline 
\texttt{QNH}& \textit{xxx hPa;} ou \textit{xxx inHg;} 
ou  \textit{xxx mmHg;}& Valeur du QNH\\
\hline 
\texttt{QFE}& \textit{xxx hPa;} ou \textit{xxx inHg;} 
ou  \textit{xxx mmHg;} & Valeur du QFE\\
\hline 
\texttt{vent}& \textit{ddd/ff kt;} ou  \textit{ddd/ff m.s-1;}
& Direction en degrés et force du vent réel\\
\hline 
\texttt{temperature}& \textit{xxx C;} ou \textit{xxx F;} & Température en degrés C ou F\\
\hline 
\texttt{point\_de\_rosee}& \textit{xxx C;} ou \textit{xxx F;} & Température en degrés C ou F\\
\hline 
\texttt{visibilite}& \textit{xxxx m;} ou \textit{xx km;}& Visibilité\\
\hline 
\texttt{humidite}& \textit{xx \%;} & Humidité relative\\
\hline
\texttt{nuage}& \textit{texte;} & Description des nuages\\
\hline
\texttt{temps\_present}& \textit{texte;} & Description du temps présent\\
\hline
\texttt{origine}& \textit{texte;} & Origine de l'information météorologique\\
\hline
\texttt{commentaire}& \textit{texte;} & Commentaire de l'opérateur\\
\hline
\end{tab3col}
\saute

Le champ commentaire sert à recueillir les remarques de
l'opérateur sur l'information météorologique. Celles-ci
peuvent par exemple, décrire l'évolution rapide d'une situation,
ou apporter des informations non fournies dans les autres champs
du paragraphe.

\subsubsection{{[}METEO\_FIN{]}}

Ce paragraphe décrit les conditions météorologiques en fin
d'acquisition.

Le format est identique à \hbox{[METEO\_DEBUT]}.


\subsubsection{{[}SEQUENCES{]}}

Ce paragraphe décrit les différentes séquences spatiales et temporelles
au cours de l'acquisition.

\saute
\noindent\begin{tab2col}
\hline 
Champ& Format valeur\\
\hline 
\texttt{palier}&
\textit{nom/hhmmss\( _{deb} \)/hmmss\( _{fin} \)/Wpt\( _{deb} \)/Wpt\( _{fin} \)/Cap/Alt./comment;}\\
\hline 
\texttt{sondage}&
\textit{nom/hhmmss\( _{deb} \)/hmmss\( _{fin} \)/Wpt\( _{deb} \)/Wpt\( _{fin} \)/Alt\( _{deb} \)/Alt\( _{fin} \)/comment;}\\
\hline 
\texttt{cercle}&
\textit{nom/hhmmss\( _{deb} \)/hmmss\( _{fin} \)/Wpt/Alt/comment;}\\
\hline 
\texttt{zoom}&
\textit{nom/hhmmss\( _{deb} \)/hmmss\( _{fin} \)/Wpt\( _{deb} \)/Wpt\( _{fin} \)/Cap/Alt\( _{deb} \)/Alt\( _{fin} \)/comment;}\\
\hline 
\texttt{virage}&
\textit{nom/hhmmss\( _{deb} \)/hmmss\( _{fin} \)/Cap\( _{av} \)/Cap\( _{ap} \)/Angle parcouru/Alt/comment;}\\
\hline 
\texttt{station}&
\textit{nom/hhmmss\( _{deb} \)/hmmss\( _{fin} \)/Wpt/Cap/Alt/comment;}\\
\hline 
\texttt{sequence}&
\textit{nom/hhmmss\( _{deb} \)/hmmss\( _{fin} \)/comment;}\\
\hline 
\end{tab2col}
\saute

 \textit{nom} : abréviation de la séquence

 \textit{hhmmss\( _{deb} \)} : \textit{}heure de début de la
séquence sous la forme hhmmss

 \textit{hmmss\( _{fin} \)} : \textit{}heure de début de la
séquence sous la forme hhmmss

 \textit{Wpt} : abréviation du point géographique où se déroule
la séquence

 \textit{Wpt\( _{deb} \)} : abréviation du point de début de la
séquence

 \textit{Wpt\( _{fin} \)} : abréviation du point de fin de la
séquence

 \textit{Cap} : cap de la plate-forme durant la séquence en
degrés

 \textit{Cap$_{av}$} : cap de la plate-forme avant la séquence en
degrés 

 \textit{Cap$_{ap}$} : cap de la plate-forme après la séquence en
degrés

 \textit{Alt.} : altitude de la plate-forme durant la séquence
suivie de l'unité ft ou m, ou notée en niveau de vol FLxxx

 \textit{Alt\( _{deb} \)} : altitude de la plate-forme au début
de la séquence suivie de l'unité ft ou m, ou notée en niveau de vol FLxxx

 \textit{Alt\( _{fin} \)} : altitude de la plate-forme au début
de la séquence suivie de l'unité ft ou m,
ou notée en niveau de vol FLxxx

 \textit{Angle parcouru} : angle parcouru par la plate-forme au cours de la
séquence. Il est noté positif dans le sens des aiguilles d'une montre.

 \textit{comment} : commentaire texte


\subsubsection{{[}ETAT\_CAPTEURS{]}}

\saute
\noindent\begin{tab3col}
\hline 
Champ& Format valeur& Description\\
\hline 
\hline 
\texttt{avant\_acq}& \textit{texte;} &
commentaire sur l'état des capteurs avant l'acquisition\\
\hline 
\texttt{pendant\_acq}& \textit{texte;} &
commentaire sur l'état des capteurs pendant l'acquisition\\
\hline 
\texttt{apres\_acq}& \textit{texte;} &
commentaire sur l'état des capteurs après l'acquisition\\
\hline 
\end{tab3col}
\saute


\subsubsection{{[}PARAMETRES\_ACQUIS{]}}

\saute
\noindent\begin{tab3col}
\hline 
Champ& Format valeur& Description valeur\\
\hline 
\hline 
\texttt{abréviation}& \textit{texte;}& nom détaillé du paramètre\\
\hline 
\end{tab3col}
\saute

Ce paragraphe contient une ligne "\texttt{champ = valeur}" par paramètre
acquis. 
Le nom du champ est l'abréviation du paramètre.\\ 
Ce paragraphe peut être très long, car pour certaines plates-formes
le nombre de paramètres est important. Par conséquent pour améliorer la lisibilité
du fichier, ce paragraphe doit être \textbf{placé en fin de fichier}.





\section{Description des campagnes de mesures}
\label{sec:campagnes}
Une campagne de mesure est un ensemble de séries de mesures provenant de une ou plusieurs
plate-formes.

\subsection{Emplacement}
\texttt{\$AIDA\_META\//campagnes}
 
\subsection{Nommage}

Le fichier est nommé du nom de la campagne officiel transformé 
en majuscules avec l'extension \texttt{.camp}.\\
Dans le cas de campagnes ayant le même nom mais à quelques années d'intervalles on rajoutera 
l'année derrière le nom.

\subsection{Format}

C'est un fichier au format "INI", divisé en paragraphes.

Il contient un paragraphe IDENTIFICATION contenant les informations générales
de la campagne.
Il contient ensuite un paragraphe par plate-forme ayant participé à la campagne, nommé par le nom de la 
plate-forme défini dans le fichier plate-formes.

\subsubsection{[IDENTIFICATION]}

\saute
\noindent\begin{tabular}{|p{4.3cm}|p{2.2cm}|p{3.9cm}|p{5.5cm}|}
\hline 
Champ& Caractère& Format valeur& Description \\
\hline 
\hline 
\texttt{nom}& \textbf{obligatoire} & \textit{texte}& Nom de la campagne, identique au nom du fichier\\
\hline 
\texttt{date\_debut}& \textbf{recommandé} & \textit{JJ/MM/AAAA} & Date de début de la campagne\\
\hline 
\texttt{date\_fin }& \textbf{recommandé} & \textit{JJ/MM/AAAA} ou \textit{running} si la  campagne n'est pas terminée & Date de fin de la campagne \\
\hline 
\texttt{projet }& \textbf{recommandé} & \textit{texte} & Acronyme du projet. Le projet est plus large que la campagne. Il peut contenir plusieurs campagnes.\\
\hline 
\texttt{description }& \textbf{obligatoire} & \textit{texte} & Description succincte en anglais de la campagne et des objectifs scientifiques. \\
\hline
\texttt{region\_geographique }& \textbf{recommandé} & \textit{texte} & Texte en anglais. Exemple : South-West of France \\
\hline
\texttt{un\_m49\_code}& \textbf{recommandé} & \textit{Code un\_m49 de la région} & Codes disponibles sous : \href{https://en.wikipedia.org/wiki/UN_M49}{https://en.wikipedia.org/wiki/
UN\_M49}.\\
\hline 
\texttt{geospatial\_lat\_min }& \textbf{recommandé} & Degrés décimaux & Latitude minimale de la zone géographique explorée au cours de la campagne\\
\hline 
\texttt{geospatial\_lat\_max }& \textbf{recommandé} & Degrés décimaux & Latitude maximale de la zone géographique explorée \\
\hline 
\texttt{geospatial\_lon\_min }& \textbf{recommandé} & Degrés décimaux & Longitude minimale de la zone géographique explorée \\
\hline 
\texttt{geospatial\_lon\_max }& \textbf{recommandé} & Degrés décimaux & Longitude maximale de la zone géographique explorée \\
\hline 
\texttt{description }& \textbf{obligatoire} & \textit{texte} & Description succincte en anglais de la campagne et des objectifs scientifiques. \\
\hline 
\texttt{principal\_investigator }& \textbf{obligatoire} & \textit{prenom1:nom1:email1:
institution1; prenom2:nom2:
email2:institution2; ...} & PI(s) de la campagne globale au même format que celui utilisé dans nos fichiers NetCDF.\\
\hline 


\end{tabular}
\saute

\subsubsection{[<nom plate-forme>]}


\saute
\noindent\begin{tabular}{|p{4.3cm}|p{2.1cm}|p{4.7cm}|p{5cm}|}
\hline 
Champ& Caractère& Format valeur& Description \\
\hline 

\texttt{date\_debut}& \textbf{recommandé} & \textit{JJ/MM/AAAA} & Date de début de des mesures de la plate-forme dans la campagne\\
\hline 
\texttt{date\_fin }& \textbf{recommandé} & \textit{JJ/MM/AAAA} ou \textit{running} si la  campagne n'est pas terminée & Date de fin de des mesures de la plate-forme dans la campagne. \\
\hline 
\texttt{principal\_investigator }& \textbf{obligatoire} & \textit{prenom1:nom1:email1:
institution1; prenom2:nom2:
email2:institution2; ...} & PI(s) de la plateforme au même format que celui utilisé dans nos fichiers NetCDF. \\
\hline 

\texttt{liste\_acqid }& \textbf{obligatoire} & Liste d'Acq\_Id séparés par des virgules.\newline
Description d'une plage d'Acq\_Id sous la forme :\newline
\textit{<Acq\_Id\_deb>$-$<Acq\_Id\_fin>}
\textit{<Acq\_Id\_fin>} peut être égal à * si la campagne est en cours & Liste des identificateurs d'acquisition (Acq\_Id) correspondant à la campagne.\\
\hline 

\end{tabular}
\saute

\section{Description des familles de paramètres }
\label{sec:familles}

Il peut être utile de regrouper les paramètres en familles, ce qui
permet des tris thématiques, des aides à la saisie, etc.
Le fichier \path{familles}
décrit la liste des familles de paramètres connues d'AIDA.

\subsection{Emplacement}
 Le fichier \path{familles} est situé dans le répertoire
 \path{${AIDA_HOME}/config/}

\subsection{Format}
   C'est un fichier au format "INI". Il est encodé en
utf-8\footnote{contrairement à la plupart des fichiers d'AIDA,
qui sont codés en ISO8859-15}.
   Il est structuré sous la forme d'un paragraphe par famille. 

Chaque paragraphe se présente comme suit :
\begin{description}
  \item[[Nom du paragraphe]] nom de la famille, en codage ASCII sans accents (par exemple, \texttt{temperature})
  \item[acronyme = ]un groupe de trois lettres
  \item[nom:fr = ] le nom en français correct, avec accents (par exemple, \texttt{température})
  \item[nom:en = ] le nom en anglais
\end{description}
  S'il apparaît un besoin de traduction dans une autre langue,
il est possible de rajouter autant de lignes \textbf{\texttt{nom:<code\_langue>}} que nécessaire,
en utilisant pour code\_langue les valeurs de la table ISO639-1 (par
exemple \texttt{nom:ca} pour le catalan)


\section{Description des aéroports}
\label{sec:aeroports}
  Certains traitements ont besoin d'informations sur des aéroports : par
exemple, l'altitude de la piste, pour calculer une altitude-pression par
intégration à partir du décollage. Ces informations sont rassemblées dans
le fichier \path{airports.fdd}

\subsection{Emplacement}
  Le fichier \path{airports.fdd} est situé dans le répertoire
  \path{${AIDA_HOME}/profils/${AIDA_PROFIL}/}

\subsection{Format}
   C'est un fichier au format "INI", contenant un paragraphe par aéroport.
Le nom du paragraphe (entre crochets) est le nom de l'aéroport,
en français, sans accents, espace ni caractère de soulignement (\texttt{\_}).
Pour les noms composés, des tirets (\texttt{-}) seront
utilisés. Exemples : \texttt{La-Corogne}, \texttt{Baden-baden},
\texttt{Nimes-Courbessac}.

  Ce nom doit correspondre à celui utilisé pour les
champs \texttt{lieu\_depart} et \texttt{lieu\_arrivee}
du paragraphe \texttt{IDENTIFICATION} du fichier .inf.

  Les attributs possibles sont :\newline
\noindent\begin{tab3col}
\hline 
Champ& Format valeur& Description valeur\\
\hline 
\hline 
\texttt{ICAO\_ID}& \textit{4 caractères majuscule}& l'indicatif OACI de l'aéroport\\
\hline 
\texttt{Z0}& \textit{valeur numérique en mètres}& L'altitude de la piste (de l'endroit où se font la plupart des décollages)\\
\hline 
\end{tab3col}
\saute

\section{Description de l'évolution des paramètres de traitement.}
\label{sec:Evolution}
  Certains traitements utilisent des paramètres qui évoluent dans le temps. Ces
paramètres sont représentés par la syntaxe \$Chrono.PARAM dans la ligne de
commande adc. La valeur de PARAM est lue dans un fichier
\texttt{nom\_plateforme.chrono.}

\subsection{Emplacement}
  Les fichiers \path{nom_plateforme.chrono} sont situés dans le répertoire
  \path{${AIDA_HOME}/profils/${AIDA_PROFIL}/chrono}
 
\subsection{Format}
 Les lignes commençant par \# sont des commentaires.
 Les lignes blanches sont autorisées.
 \\
 
 Le fichier se compose d'une section par paramètre. La section est identifiée par le nom du paramètre entre crochets : [PARAM].
 Chaque section contient des lignes qui décrivent une plage de dates et la valeur de PARAM à l'intérieur de cette plage. Les bornes de la plage sont comprises. Le séparateur est constitué d'un ou plusieurs espaces ou tabulations.
 \\
 
 \textbf{Format des dates} : YYYY-MM-DDTHH:MI:SSZ
 
 \textbf{Valeurs possibles de PARAM} :
 
 -nombre (ex: -1.5)
        
 -chaîne (ex: "celsius")
      
 -nom de vecteur (ex: hum\_sol3cm\_ec5\_c3\_\%900)
 \\
 
 
 \textbf{Exemple de fichier chrono} : 
 
 \begin{lstlisting}
 # Chrono toulouse-flux
 
 # délai du LICOR
 [DELAI]
 2012-06-15T00:00:00Z    2014-10-14T23:59:59Z    -0.200
 2014-10-15T00:00:00Z    2015-05-28T23:59:59Z    -0.250
 2015-05-29T00:00:00Z    2020-12-31T23:59:59Z    -0.200
 
 # unité températures pour tests
 [UNITE_TPR]
 2012-06-15T00:00:00Z    2014-07-09T23:59:59Z    "Kelvin"
 2014-07-10T00:00:00Z    2020-12-31T23:59:59Z    "Celsius"
 
 # vecteur pour tests
 [VECTEUR]
 2014-06-15T00:00:00Z    2014-07-09T23:59:59Z    hum_sol3cm_ec5_c3_%900

 2014-07-10T00:00:00Z    2020-12-31T23:59:59Z    hum_sol3cm_ec5_c3_bis_%900
 \end{lstlisting}

 \saute 
 
 
